\BeleskeID{09_2-s051}
% element: 09_2-s051:e04
% element: 09_2-s051:e05
Parazitne kapacitivnosti u odnosu na izlazni signal i baze tranzistora imaju diferencijatorski efekat. Što je veća brzina promene izlaznog napona, izraženiji je pobudni impuls. Pozitivan impuls na bazi $T_N$ može biti dovoljan da ga uključi.

Kada $T_N$ provede, izaziva pad napona na $R_{W1}$, što dovodi do uključenja $T_P$. Uključenje tranzistora $T_P$ izaziva pojavu napona na $R_{S1}$ i time bolje provođenje $T_N$. To zatim povećava pad na $R_{W1}$, dodatno uključuje $T_P$, povećava napon na $R_{S1}$ i ponovo pojačava provođenje $T_N$. Pozitivna reakcija traje dok oba tranzistora ne ostanu uključena bez dodatne spoljašnje pobude. Time se ostvaruje putanja veoma velike struje između napajanja i mase.

Ista reakcija može početi uključenjem $T_P$ usled negativnog impulsa iza diferencijatora. Redosled početne pobude je različit, ali se uspostavlja ista regenerativna petlja.
